\exercice{Oscillateur à filtre RLC}
\enonce{
On considère le montage suivant, constitué d'un montage amplificateur non-inverseur et d'un filtre RLC-série.
\begin{center}
  \begin{circuitikz}
    \draw (0,0) node[en amp, yscale=-1] (ALI) {};
    \draw (ALI.-) --++ (0,-1.7) to[R=$R_1$] ++ (0,-2) node[ground] {};
    \draw (ALI.-) ++ (0,-1.7) to[R=$R_2$] ++ (2.2,0) -| (ALI.out) to[L=$L$] ++ (2,0) to[R=$r$] ++ (2,0) to[C=$C$] ++ (2,0) coordinate (sortie) to[R=$R$] ++ (0,-4.035) node[ground] {};
    \draw (sortie) --++ (0,2) -| (ALI.+);
    \draw[<-,>=latex] (sortie) ++ (1,0) --++ (0,-4.035) node[midway, right] {$s$};
    \draw[<-,>=latex] (ALI.out) ++ (0.2,-0.1) --++ (0,-3.935) node[midway, right] {$e$};
  \end{circuitikz}
\end{center}}
\begin{enumerate}
    \question{Déterminer la fonction de transfert $H(p)=\dfrac{S(p)}{E(p)}$ du filtre RLC-série.}
      [Déterminer l'impédance équivalente de $L$, $r$ et $C$ en série puis utiliser un pont diviseur de tension.]
    \question{Donner la fonction de transfert du montage amplificateur non-inverseur.}
    \question{À quelle condition observe-t-on des oscillations quasi-sinusoïdales ?}
      [Utiliser les deux fonctions de transfert des questions précédentes.]
      [Remplacer $E(p)$ dans la fonction de transfert du filtre en utilisant la fonction de transfer de l'amplificateur non-inverseur puis éliminer $E(p)$]
      [Éliminer les traits de fraction puis prendre la partie réelle et la partie imaginaire.]
    \question{Quelle est l'amplitude de $e(t)$ ? Quelle est celle de $s(t)$ ?}
      [Qu'est-ce qui limite l'amplitude des oscillations ?]
      [Quelle est l'amplitude maximale de la tension de sortie d'un ALI ?]
      [Utiliser une des fonctions de transfert pour relier les amplitudes de $s$ et de $e$.]
    \question{Laquelle de ces deux tensions est la << plus sinusoïdale >> ?}
      [Est-ce le filtre passe bande ou l'ALI qui a tendance à "purifier" le spectre ?]
    \question{À quelle condition les oscillations démarrent-elles ?}
      [Établir une équation différentielle portant sur $e$ ou $s$.]
      [Combiner les fonctions de transfert pour éliminer $E$ ou $S$, supprimer les traits de fraction puis passer en temporel.]
\end{enumerate}

\exercice{Oscillateur à décharge de condensateur}
\enonce{
On étudie le circuit suivant, où $R_1=\SI{100}{k\ohm}$, $R_2=\SI{10}{k\ohm}$, $R_3=\SI{10}{k\ohm}$ et $C=\SI{10}{nF}$.
\begin{center}
  \begin{circuitikz}
    \draw (0,0) node[en amp] (ALI) {};
    \draw (ALI.+) --++ (0,-1.7) to[R=$R_2$] ++ (0,-2) node[ground] {};
    \draw (ALI.+) --++ (0,-1.7) to[R=$R_1$] ++ (2.2,0) -| (ALI.out);
    \draw (ALI.out) node[above] {$V_s$} to[R=$R_3$,*-*] ++ (2,0) coordinate (sortie) to[C=$C$] ++ (0,-4.035) node[ground] {};
    \draw (sortie) node[right] {$V_-$} --++ (0,2) -| (ALI.-);
  \end{circuitikz}
\end{center}}
\begin{enumerate}
    \question{Identifier le filtre passif inclus dans ce montage. Quelle est sa fonction de transfert ? Quelle est l'équation différentielle associée ?}
      [Le montage est constitué d'un comparateur à hystérésis et d'un filtre qu'il s'agit d'identifier.]
      [Le filtre RC-série est-il un filtre passe-haut ou passe-bas ? Quelle est la forme canonique d'une fonction de transfer de ce type ?]
    \question{Résoudre cette équation différentielle en supposant $V_s$ constant.}
    \question{Quel montage de l'ALI reconnait-on dans ce montage ? Donner sa caractéristique $(V_s,V_-)$.}
      [En cas d'hésitation entre comparateur à hystérésis positif et négatif, regarder où se fait l'entrée.]
    \question{En supposant qu'à $t=0$, $V_s=+V_\text{sat}$ et $V_-=$, tracer $V_-(t)$ et $V_s(t)$.}
      [Tracer $V_-(t)$ en supposant que $V_s(t)=+V_\text{sat}$. Jusqu'à quand ce tracé reste-t-il valable ?]
      [À quelle condition $V_s$ passe-t-il de $+V_\text{sat}$ à $-V_\text{sat}$ ?]
    \question{Que vaut la période des signaux produits ?}
      [Déterminer la demi-période, c'est-à-dire le temps nécessaire pour que $V_-$ passe de $\frac{R_2}{R_1+R_2}V_\text{sat}$ à $\frac{R_2}{R_1+R_2}V_\text{sat}$.]
      [Résoudre complètement (constante comprise) l'équation différentielle vérifiée par $V_-$ sur une demi période. On pourra appeler $t_1$ le début de cette demi-période et $t_2$ sa fin.]
\end{enumerate}

\exercice{Oscillateur à résistance négative}
\enonce{
On considère le montage ci-dessous, appelé montage à résistance négative. On \textbf{suppose} l'ALI en fonctionnement stable.
\begin{center}
  \begin{circuitikz}
    \draw (0,0) node[en amp, yscale=-1] (ALI) {};
    \draw (ALI.-) --++ (0,-1.7) to[R=$R_1$] ++ (0,-2) node[ground] {};
    \draw (ALI.-) ++ (0,-1.7) to[R=$R_1$] ++ (2.2,0) -| (ALI.out);
    \draw (ALI.+) --++ (0,1.7) to[R=$R$] ++ (2.2,0) -| (ALI.out);
    \draw (ALI.+) to[short, i<=$i$] ++ (-1,0) coordinate (entree);
    \draw[<-,>=latex] (entree) ++ (0,-0.1) --++ (0,-4.26) coordinate (masse) node[midway, left] {$u$};
    \draw (masse) node[ground] {};
  \end{circuitikz}
\end{center}}
\begin{enumerate}
    \question{Déterminer une relation entre la tension $u$ et la tension de sortie de l'ALI.}
      [Le montage est-il stable ? Que peut-on dire de l'entrée différentielle $\epsilon$ ?]
      [À l'aide d'un pont diviseur de tension entre les résistances $R_1$ et d'une loi des mailles, relier l'entrée différentielle $\epsilon$ à $u$ et la sortie de l'ALI.]
    \question{En utilisant la loi d'Ohm, en déduire l'impédance d'entrée du montage $Z_e=\frac{\underline{u}}{\underline{i}}$.}
      [Écrire la loi d'Ohm dans la résistance $R$.]
      [Utiliser la fonction de transfer de la question précédente pour éliminer la tension de sortie de l'ALI dans la loi d'Ohm.]
\end{enumerate}

\enonce{
Ce montage, qui se comporte comme une <<résistance négative>>, est placée dans le circuit suivant où $R_1=\SI{10}{k\ohm}$ et $r=\SI{10}{k\ohm}$.
\begin{center}
  \begin{circuitikz}
    \draw (0,0) to[short, i<=$i$] ++ (-1,0) coordinate (entree) to[R, l_=$r$] ++ (-2,0) to[C, l_=$C$] ++ (0,-2) to[L, l_=$L$] ++ (0,-2) --++ (3,0) to[R,l=$Z_e$] (0,0);
    \draw[<-,>=latex] (entree) ++ (0,-0.1) --++ (0,-3.8) node[midway, left] {$u$};
  \end{circuitikz}
\end{center}
}
\begin{enumerate}
  \setcounter{enumi}{2}
  \question{Déterminer une équation différentielle sur $i$.}
    [Appliquer la loi des mailles et utiliser les caractéristiques de différents composants.]
    [Dériver la loi des mailles et utiliser que $u_L=-L\frac{di}{dt}$, $i_c=-C\frac{di}{dt}$, $u_r=-Ri$ et $u=Z_ei$.]
  \question{Sous quelles conditions sur la valeur de $R$ les oscillations démarrent-elles ?}
    [A quelle condition les solutions de l'équation différentielle précédente divergent-elles ?]
\end{enumerate}

\exercice{Hartley oscillator}
\enonce{
First, we will study the electronic filter below, named Hartley filter.
\begin{center}
  \begin{circuitikz}
    \draw (0,0) to[R=$R$, -*] (2,0) node[above] {$A$} to[C=$C$] (2,-2) node[ground] {};
    \draw (2,0) to[L=$L$] (4,0) to[L=$L$] (4,-2) node[ground] {};
    \draw[<-, >=latex] (0,-0.1) -- (0,-2) node[midway, left] {$e$};
    \draw (0,-2) node[ground] {};
    \draw[<-, >=latex] (5,0) -- (5,-2) node[midway, right] {$s$};
  \end{circuitikz}
\end{center}}
\begin{enumerate}
    \question{Using Kirchhoff's nodal rule in $A$, express its potential as a function of the voltages $e$ and $s$.}
      [Si la relation des nœuds en tension n'est pas connue, appliquer la loi des nœuds et remplacer chaque courant en faisant apparaitre une tension grâce à la loi d'Ohm (en complexes) dans $R$, $L$ et $C$.]
    \question{Ascertain $s$ as a function of the potential in $A$ thanks to the voltage diviser law. With the help of the previous question, determine the transfer function of the Hartley filter.}
\end{enumerate}

\enonce{
The complete Hartley oscillator's electrical schema is shown below.
\begin{center}
  \begin{circuitikz}
    \draw (0,0) node[en amp, yscale=-1] (ALI) {};
    \draw (ALI.-) --++ (0,-1.7) to[R=$R_1$] ++ (0,-2) node[ground] {};
    \draw (ALI.-) --++ (0,-1.7) to[R=$R_2$] ++ (2.2,0) -| (ALI.out);
    \draw (ALI.out) to[R=$R$, -*] ++ (2,0) to[C=$C$] ++ (0,-4.035) node[ground] {};
    \draw (ALI.out) ++ (2,0) to[L=$L$] ++ (2,0) coordinate (sortie) to[L=$L$] ++ (0,-4.035) node[ground] {};
    \draw[<-,>=latex] (sortie) ++ (1,0) --++ (0,-4) node[midway, right] {$s$};
    \draw[<-,>=latex] (ALI.out) ++ (0.2,-0.1) --++ (0,-4) node[midway, right] {$e$};
    \draw (sortie) --++ (0,2) -| (ALI.+);
  \end{circuitikz}
\end{center}}
\begin{enumerate}\setcounter{enumi}{2}
    \question{Which operational amplifier assembly can be recognized ? Give its transfer function without any demonstration.}
      [Pour reconnaitre le montage à ALI, celui-ci est-il stable ou instable ? L'entrée est-elle "côté" borne inverseuse ou  non-inverseuse ?]
    \question{Under which condition do the oscillations \textbf{start} ?}
      [Obtenir une équation différentielle portant sur $s$ ou sur $e$.]
      [Transformer la fonction de transfert du filtre de Hartley en une équation différentielle sur $s$ et $e$. Remplacer l'un des deux en utilisant la relation entrée-sortie d'un amplificateur non-inverseur.]
      [Les solutions de l'équation différentielle doivent-elles converger ou diverger pour que les oscillations démarrent.]
    \question{How sinusoidal oscillations can be obtained ? What will be their frequency ?}
      [En partant de la fonction de transfert du filtre de Hartley, éliminer $E$ et $S$.]
      [Éliminer les fractions de l'équation ainsi obtenue et en prendre partie réelle et partie imaginaire.]
    \question{What is the amplitude of $e$ ? What is the one of $s$ ?}
      [Quelles sont les amplitudes de $e$ et $s$ ?]
      [Qu'est-ce qui limite l'amplitude des oscillations ?]
    \question{Which voltage will be the most sinusoidal ?}
      [Est-ce l'amplificateur non-inverseur ou le filtre de Hartley qui "purifie" le spectre de son entrée ?]
\end{enumerate}